% ============================================================
% setup/format.tex —— 版面与排版规范（按目标期刊微调）
% 两刊要求对照（详见 README.md）：
%   《系统工程与电子技术》：图表题/注释中文；坐标轴标量和单位；
%     物理量 GB 3100～3102-1993（变量斜体，矢量矩阵黑斜体，
%     集合黑正体，常数/运算符/函数正体）
%   《信息与控制》：图题表题中英双语；公式编号置右；
%     投稿为 Word 双栏（排版终稿由编辑部完成，本模板用于写作）
% 注：各刊细节不同，投稿前以目标期刊"稿件模板/征稿简则"为准。
% ============================================================

% ---------- 页边距（期刊通行：上下 2.5cm，左右 2.0~2.5cm） ----------
\geometry{top=2.5cm, bottom=2.5cm, left=2.2cm, right=2.2cm}

% ---------- 各级标题（黑体；标题层级= 节 / 小节） ----------
\ctexset{
  section = {
    format     = \heiti\zihao{4},      % 节标题：黑体四号
    beforeskip = 1.5ex plus .2ex minus .2ex,
    afterskip  = 1.2ex plus .2ex minus .2ex,
    indent     = 0pt
  },
  subsection = {
    format     = \heiti\zihao{-4},     % 小节标题：黑体小四
    beforeskip = 1.2ex plus .2ex minus .2ex,
    afterskip  = 1ex plus .2ex minus .2ex,
    indent     = 0pt
  },
  subsubsection = {
    format     = \heiti\zihao{5},      % 小小节标题：黑体五号
    beforeskip = 1ex plus .2ex minus .2ex,
    afterskip  = 1ex plus .2ex minus .2ex,
    indent     = 0pt
  }
}

% ---------- 图表题（五号；图题在下、表题在上） ----------
% 《系统工程与电子技术》：图表题和注释用中文
% 《信息与控制》：图题/表题中英双语（\caption{中文题}\caption*{English title}
%   或改用 bicaption 宏包）；投稿时按目标刊切换
\captionsetup{font=small, labelsep=quad, skip=4pt}
\captionsetup[figure]{position=bottom}
\captionsetup[table]{position=top}

% ---------- 公式编号：按节编号 (1)、(2)… ----------
\numberwithin{equation}{section}

% ---------- 段落 ----------
% 中文首行缩进 2 字符（ctex 默认），行距保持紧凑（期刊默认单倍行距）
\linespread{1.1}
\setlength{\parindent}{2em}
